\documentclass[10pt,a4paper]{amsart}
\usepackage{amsmath,amssymb,amsthm,graphicx}
% Retained minimal preamble for the Stacks Project chapter more-morphisms.tex (More on Morphisms).
% Verbatim excerpt of the chapter preamble: Xy-pic setup (native preamble lines 14-18) and the native
% theorem declarations with the mathematical macros (native preamble lines 154-217). Copied mechanically;
% nothing added, removed or redefined. The wrapper class and the amsmath/amssymb/amsthm/graphicx base
% packages are supplied by the assembler, so the class/externaldocument part of the native preamble is omitted.
% For commutative diagrams we use Xy-pic
\usepackage[all]{xy}

% We use 2cell for 2-commutative diagrams.
\xyoption{2cell}
% Theorem environments.
%
\theoremstyle{plain}
\newtheorem{theorem}[subsection]{Theorem}
\newtheorem{proposition}[subsection]{Proposition}
\newtheorem{lemma}[subsection]{Lemma}

\theoremstyle{definition}
\newtheorem{definition}[subsection]{Definition}
\newtheorem{example}[subsection]{Example}
\newtheorem{exercise}[subsection]{Exercise}
\newtheorem{situation}[subsection]{Situation}

\theoremstyle{remark}
\newtheorem{remark}[subsection]{Remark}
\newtheorem{remarks}[subsection]{Remarks}

\numberwithin{equation}{subsection}

% Macros
%
\def\lim{\mathop{\mathrm{lim}}\nolimits}
\def\colim{\mathop{\mathrm{colim}}\nolimits}
\def\Spec{\mathop{\mathrm{Spec}}}
\def\Hom{\mathop{\mathrm{Hom}}\nolimits}
\def\Ext{\mathop{\mathrm{Ext}}\nolimits}
\def\SheafHom{\mathop{\mathcal{H}\!\mathit{om}}\nolimits}
\def\SheafExt{\mathop{\mathcal{E}\!\mathit{xt}}\nolimits}
\def\Sch{\mathit{Sch}}
\def\Mor{\mathop{\mathrm{Mor}}\nolimits}
\def\Ob{\mathop{\mathrm{Ob}}\nolimits}
\def\Sh{\mathop{\mathit{Sh}}\nolimits}
\def\NL{\mathop{N\!L}\nolimits}
\def\CH{\mathop{\mathrm{CH}}\nolimits}
\def\proetale{{pro\text{-}\acute{e}tale}}
\def\etale{{\acute{e}tale}}
\def\QCoh{\mathit{QCoh}}
\def\Ker{\mathop{\mathrm{Ker}}}
\def\Im{\mathop{\mathrm{Im}}}
\def\Coker{\mathop{\mathrm{Coker}}}
\def\Coim{\mathop{\mathrm{Coim}}}

% Boxtimes
%
\DeclareMathSymbol{\boxtimes}{\mathbin}{AMSa}{"02}

%
% Macros for moduli stacks/spaces
%
\def\QCohstack{\mathcal{QC}\!\mathit{oh}}
\def\Cohstack{\mathcal{C}\!\mathit{oh}}
\def\Spacesstack{\mathcal{S}\!\mathit{paces}}
\def\Quotfunctor{\mathrm{Quot}}
\def\Hilbfunctor{\mathrm{Hilb}}
\def\Curvesstack{\mathcal{C}\!\mathit{urves}}
\def\Polarizedstack{\mathcal{P}\!\mathit{olarized}}
\def\Complexesstack{\mathcal{C}\!\mathit{omplexes}}
% \Pic is the operator that assigns to X its picard group, usage \Pic(X)
% \Picardstack_{X/B} denotes the Picard stack of X over B
% \Picardfunctor_{X/B} denotes the Picard functor of X over B
\def\Pic{\mathop{\mathrm{Pic}}\nolimits}
\def\Picardstack{\mathcal{P}\!\mathit{ic}}
\def\Picardfunctor{\mathrm{Pic}}
\def\Deformationcategory{\mathcal{D}\!\mathit{ef}}
\newenvironment{reference}{}{}
\newenvironment{slogan}{}{}
\newenvironment{history}{}{}
\usepackage{cleveref}
\crefname{originallemma}{original lemma}{original lemmas}
\Crefname{originallemma}{Original lemma}{Original lemmas}
\makeatletter
\newlabel{algebra-lemma-formally-etale-etale}{{00UR}{source}{}{sourcepointer-1}{}}
\newlabel{algebra-lemma-formally-etale-etale@cref}{{[originallemma][0][]00UR}{[1][0][]source}{}{}{}}
\makeatother
\crefname{originallemma}{original lemma}{original lemmas}
\Crefname{originallemma}{Original lemma}{Original lemmas}
\makeatletter
\newlabel{lemma-affine-formally-etale}{{lemma-affine-formally-etale}{source}{}{sourcepointer-2}{}}
\newlabel{lemma-affine-formally-etale@cref}{{[originallemma][0][]lemma-affine-formally-etale}{[1][0][]source}{}{}{}}
\makeatother
\crefname{originallemma}{original lemma}{original lemmas}
\Crefname{originallemma}{Original lemma}{Original lemmas}
\makeatletter
\newlabel{lemma-formally-etale-on-opens}{{lemma-formally-etale-on-opens}{source}{}{sourcepointer-3}{}}
\newlabel{lemma-formally-etale-on-opens@cref}{{[originallemma][0][]lemma-formally-etale-on-opens}{[1][0][]source}{}{}{}}
\makeatother
\crefname{originallemma}{original lemma}{original lemmas}
\Crefname{originallemma}{Original lemma}{Original lemmas}
\makeatletter
\newlabel{lemma-unramified-flat-formally-etale}{{lemma-unramified-flat-formally-etale}{source}{}{sourcepointer-4}{}}
\newlabel{lemma-unramified-flat-formally-etale@cref}{{[originallemma][0][]lemma-unramified-flat-formally-etale}{[1][0][]source}{}{}{}}
\makeatother
\crefname{originallemma}{original lemma}{original lemmas}
\Crefname{originallemma}{Original lemma}{Original lemmas}
\makeatletter
\newlabel{morphisms-lemma-etale-characterize}{{02GJ}{source}{}{sourcepointer-5}{}}
\newlabel{morphisms-lemma-etale-characterize@cref}{{[originallemma][0][]02GJ}{[1][0][]source}{}{}{}}
\makeatother
\crefname{originallemma}{original lemma}{original lemmas}
\Crefname{originallemma}{Original lemma}{Original lemmas}
\makeatletter
\newlabel{morphisms-lemma-locally-finite-presentation-characterize}{{01TQ}{source}{}{sourcepointer-6}{}}
\newlabel{morphisms-lemma-locally-finite-presentation-characterize@cref}{{[originallemma][0][]01TQ}{[1][0][]source}{}{}{}}
\makeatother
\crefname{originallemma}{original lemma}{original lemmas}
\Crefname{originallemma}{Original lemma}{Original lemmas}
\makeatletter
\newlabel{morphisms-section-etale}{{02GH}{source}{}{sourcepointer-7}{}}
\newlabel{morphisms-section-etale@cref}{{[originallemma][0][]02GH}{[1][0][]source}{}{}{}}
\makeatother
\title{Etale morphisms of schemes are exactly the locally finitely presented formally etale ones}
\author{The Stacks Project Authors}
\date{stacks-project:more-morphisms.tex:local-snapshot-20261011; GNU FDL 1.2 or later}
\begin{document}
\maketitle
\noindent\textit{Editorial note.} Statement span is lines 1705-1714 of sources/stacks-more-morphisms/pinned-original-source.tex; proof span is lines 1716-1742, the complete author proof (25 body lines).\par\medskip
\noindent\textit{Editorial note.} Navigation stubs: 7 cited-but-not-included labels are supplied as non-load-bearing pointers with their pinned tags where the snapshot tag index has them.\par\medskip
\noindent\textit{Editorial note.} extra\_packages is empty on purpose: the retained preamble already loads xy with options and a second declaration clashes in the assembler wrapper.\par\medskip
\noindent\textit{Editorial note.} Text is retained verbatim, including the author's own typos and placeholders.\par\medskip
\noindent\textit{Allowed standard prerequisites (editorial).}\begin{itemize}
\item Explicit prerequisite boundary. Formally etale morphisms, the decomposition of etaleness into flat, unramified and locally finitely presented, and the algebra characterisation of formally etale ring maps; the converse is retained.
\end{itemize}\medskip
\ifdefined\hypertarget\hypertarget{sourcepointer-1}{}\fi
\noindent\textit{Source reference.} Original lemma 00UR: sources/stacks-more-morphisms/pinned-original-source.tex, source lines 1742--1742. Cited by the retained statement or proof but not part of the collected spans; supplied only as a navigation pointer. The argument does not depend on its content.\par\medskip
\ifdefined\hypertarget\hypertarget{sourcepointer-2}{}\fi
\noindent\textit{Source reference.} Original lemma lemma-affine-formally-etale: sources/stacks-more-morphisms/pinned-original-source.tex, source lines 1655--1661. Cited by the retained statement or proof but not part of the collected spans; supplied only as a navigation pointer. The argument does not depend on its content.\par\medskip
\ifdefined\hypertarget\hypertarget{sourcepointer-3}{}\fi
\noindent\textit{Source reference.} Original lemma lemma-formally-etale-on-opens: sources/stacks-more-morphisms/pinned-original-source.tex, source lines 1574--1578. Cited by the retained statement or proof but not part of the collected spans; supplied only as a navigation pointer. The argument does not depend on its content.\par\medskip
\ifdefined\hypertarget\hypertarget{sourcepointer-4}{}\fi
\noindent\textit{Source reference.} Original lemma lemma-unramified-flat-formally-etale: sources/stacks-more-morphisms/pinned-original-source.tex, source lines 1626--1628. Cited by the retained statement or proof but not part of the collected spans; supplied only as a navigation pointer. The argument does not depend on its content.\par\medskip
\ifdefined\hypertarget\hypertarget{sourcepointer-5}{}\fi
\noindent\textit{Source reference.} Original lemma 02GJ: sources/stacks-more-morphisms/pinned-original-source.tex, source lines 1742--1742. Cited by the retained statement or proof but not part of the collected spans; supplied only as a navigation pointer. The argument does not depend on its content.\par\medskip
\ifdefined\hypertarget\hypertarget{sourcepointer-6}{}\fi
\noindent\textit{Source reference.} Original lemma 01TQ: sources/stacks-more-morphisms/pinned-original-source.tex, source lines 1742--1742. Cited by the retained statement or proof but not part of the collected spans; supplied only as a navigation pointer. The argument does not depend on its content.\par\medskip
\ifdefined\hypertarget\hypertarget{sourcepointer-7}{}\fi
\noindent\textit{Source reference.} Original lemma 02GH: sources/stacks-more-morphisms/pinned-original-source.tex, source lines 1742--1742. Cited by the retained statement or proof but not part of the collected spans; supplied only as a navigation pointer. The argument does not depend on its content.\par\medskip
\begin{lemma}
\label{lemma-etale-formally-etale}
Let $f : X \to S$ be a morphism of schemes.
The following are equivalent:
\begin{enumerate}
\item The morphism $f$ is \'etale, and
\item the morphism $f$ is locally of finite presentation and
formally \'etale.
\end{enumerate}
\end{lemma}

\begin{proof}
Assume $f$ is \'etale.
An \'etale morphism is locally of finite presentation, flat and unramified, see
Morphisms, Section \ref{morphisms-section-etale}.
Hence $f$ is locally of finite presentation and formally \'etale, see
Lemma \ref{lemma-unramified-flat-formally-etale}.

\medskip\noindent
Conversely, suppose that $f$ is locally of finite presentation and
formally \'etale. Being \'etale is local in the Zariski topology on
$X$ and $S$, see
Morphisms, Lemma \ref{morphisms-lemma-etale-characterize}.
By
Lemma \ref{lemma-formally-etale-on-opens}
we can cover $X$ by affine opens $U$ which map into affine opens
$V$ such that $U \to V$ is formally \'etale (and of finite presentation, see
Morphisms,
Lemma \ref{morphisms-lemma-locally-finite-presentation-characterize}).
By
Lemma \ref{lemma-affine-formally-etale}
we see that the ring maps $\mathcal{O}(V) \to \mathcal{O}(U)$ are
formally \'etale (and of finite presentation).
We win by
Algebra, Lemma \ref{algebra-lemma-formally-etale-etale}.
(We will give another proof of this implication when we discuss
formally smooth morphisms.)
\end{proof}

\end{document}
